\section{Тема 2. Разветвляющиеся алгоритмы}

\subsection{Постановка задачи (вариант~4)}

Четырёхугольник $ABCD$ задан координатами своих вершин на плоскости:
$A(x_1,y_1)$, $B(x_2,y_2)$, $C(x_3,y_3)$, $D(x_4,y_4)$.
Определить тип четырёхугольника: прямоугольник, параллелограмм,
трапеция, произвольный четырёхугольник.

\subsection{Математическое обоснование}

Две стороны $AB$ и $CD$ параллельны, если векторное произведение
направляющих векторов равно нулю:
\[
  \vec{AB} \times \vec{DC} = 0
  \;\Leftrightarrow\;
  (x_2{-}x_1)(y_4{-}y_3) - (y_2{-}y_1)(x_4{-}x_3) = 0.
\]
Аналогично проверяется параллельность $BC$ и $AD$.
Если обе пары противоположных сторон параллельны~--- это
параллелограмм; если параллельна ровно одна пара~--- трапеция;
если ни одна~--- произвольный четырёхугольник.

Параллелограмм является прямоугольником, если смежные стороны
перпендикулярны:
\[
  \vec{AB} \cdot \vec{AD} = 0.
\]

\subsection{Блок-схема алгоритма}

\begin{center}
\begin{tikzpicture}[
  node distance=10mm, start chain=going below,
  nodes={draw=black, top color=white, bottom color=lightgray!50,
         minimum width=5.4cm, align=center, font=\small},
  arrow/.style={->, >=Stealth, thick},
]
  \node[draw, rounded corners, on chain] {Start};
  \node[shape=trapezium, trapezium left angle=70,
        trapezium right angle=110, on chain, draw, join=by arrow]
       {$A,B,C,D$};
  \node[shape=rectangle, on chain, draw, join=by arrow]
       {$p_1 := \mathrm{par}(AB,DC)$,\quad $p_2 := \mathrm{par}(BC,AD)$};
  \node[shape=diamond, aspect=2, on chain, draw, join=by arrow,
        minimum width=5.6cm]
       (c1) {$p_1 \wedge p_2$\,?};
  \node[shape=diamond, aspect=2, draw, right=2.4cm of c1,
        minimum width=4.6cm]
       (c2) {$AB \perp AD$\,?};
  \draw[arrow] (c1.east) -- node[above,font=\scriptsize]{$+$} (c2.west);
  \node[shape=trapezium, trapezium left angle=70,
        trapezium right angle=110, draw, above right=5mm and 5mm of c2]
       (rect) {«Прямоугольник»};
  \node[shape=trapezium, trapezium left angle=70,
        trapezium right angle=110, draw, below right=5mm and 5mm of c2]
       (par) {«Параллелограмм»};
  \draw[arrow] (c2.north east) -- node[above,font=\scriptsize]{$+$} (rect.west);
  \draw[arrow] (c2.south east) -- (par.west);
  \node[shape=diamond, aspect=2, on chain, draw,
        below=16mm of c1, minimum width=5.6cm]
       (c3) {$p_1 \vee p_2$\,?};
  \draw[arrow] (c1.south) -- (c3.north);
  \node[shape=trapezium, trapezium left angle=70,
        trapezium right angle=110, draw, right=2.4cm of c3]
       (trap) {«Трапеция»};
  \draw[arrow] (c3.east) -- node[above,font=\scriptsize]{$+$} (trap.west);
  \node[shape=trapezium, trapezium left angle=70,
        trapezium right angle=110, on chain, draw, below=12mm of c3]
       (arb) {«Произвольный»};
  \draw[arrow] (c3.south) -- (arb.north);
  \node[draw, rounded corners, on chain, draw, below=10mm of arb]
       (end) {End};
  \draw[arrow] (arb.south) -- (end.north);
  \draw[arrow] (rect.south) |- (end.east);
  \draw[arrow] (par.south)  |- (end.east);
  \draw[arrow] (trap.south) |- (end.east);
\end{tikzpicture}
\end{center}

\subsection{Текст программы}

\lstinputlisting[style=Cstyle]{../src/t02.c}

\subsection{Тестовые данные}

\begin{center}
\begin{tabular}{l|l}
\toprule
Вершины $A,B,C,D$ & Результат \\
\midrule
$(0,0),(4,0),(4,3),(0,3)$ & Прямоугольник \\
$(0,0),(4,0),(5,3),(1,3)$ & Параллелограмм \\
$(0,0),(4,0),(3,2),(1,2)$ & Трапеция \\
$(0,0),(4,0),(3,2),(0,3)$ & Произвольный четырёхугольник \\
\bottomrule
\end{tabular}
\end{center}

\textbf{Результат выполнения программы} (первый тест):
\begin{lstlisting}[language={},basicstyle=\ttfamily\small,frame=single]
Введите координаты A B C D (x y для каждой):
0 0  4 0  4 3  0 3
Прямоугольник
\end{lstlisting}
